% page 348 of the facsimile (image p-347 of the scan)
% block 152.4 x 242.0 mm ; left margin 8.0 mm
\pgc{-12.700mm}{0.000mm}{
\l{\cel{3}340}
\l{}
\l{}
\l{\sou{lingu}o.- Ta maniero expresar sua penso per parolo o skribo,}
\l{\cel{3}propra a ta o ca naciono. DEFIRS.}
\l{}
\l{\sou{linguistiko}. La cienco di la persono qua su konsakras a la}
\l{\cel{3}studio di la lingui. - DEFIRS.}
\l{}
\l{\sou{linimento}. (\sou{farmaci}o). Medikamento grasatra, uncioniva, por}
\l{\cel{3}la fricioni. - DEFIS.}
\l{}
\l{\sou{linjo}.\cel{2}Tuko ek filo o kotono, specaligita por diversa}
\l{\cel{3}uzi hemala. - F.}
\l{}
\l{\sou{\textquotedbl{}linoleum\textquotedbl{}}. Sorto di texajo nepermeebla, ek telo de juto,}
\l{\cel{3}indutita ye lin-oleo e ye pudro de korko, uzata kom}
\l{\cel{3}tapiso-materio. - DEFIRS.}
\l{}
\l{\sou{linono}. Telo de lino, tre klara, kun apreto ferma, por robi,}
\l{\cel{3}shultro-tuki, mulier-kamizi. - DFIS.}
\l{}
\l{\sou{lintelo}. Traverso ye qua konsistas la parto supra di pordo,}
\l{\cel{3}di fenestro kun aperturo rektangula, e subtenanta la}
\l{\cel{3}masonuro. - EFS.}
\l{}
\l{\sou{lipitudo}. (\sou{patol.}) La stando di persono di qua la okuli}
\l{\cel{3}esas seboza. -}
\l{}
\l{\sou{lipomo}. (\sou{patol.}) Tumoro grasoza. - DEFIRS.}
\l{}
\l{\sou{lipotimio}. (\sou{patol.})\cel{2}Esvano.}
\l{}
\l{\sou{liquacar}. (\sou{teknol.)}\cel{2}Separar un ek la metali kontenata en}
\l{\cel{3}erco, en aloyuro. - EFIS.}
\l{}
\l{\sou{liquid}a. Di qua la molekuli, poke interadheranta, fluas,}
\l{\cel{3}obediante separe la legi gravitala kande li ne konten-}
\l{\cel{3}esas da recipiento. - EFIS.}
\l{}
\l{\sou{liquidacar.}\cel{2}(\sou{trans.})\cel{2}Konkluzar estado financala per pagar}
\l{\cel{3}la pasivo e per atribuar a la yurozo la aktivo rest-}
\l{\cel{3}anta. - DEFIRS.}
\l{}
\l{\sou{liquoro}. I. Drinkajo alkoholoza, obtenata per distilo. -}
\l{\cel{3}II. Mixajo liquida uzata farmaciale, industriale, e c.}
\l{\cel{3}- DEFIRS.}
\l{}
\l{\sou{liro}. I. Muzik-instrumento kordoza, quan uzis la antiqui.}
\l{\cel{3}- II. (\sou{astron.}) Stelaro en la hemisfero nordala. - III.}
\l{\cel{3}(\sou{zool.}) Ucelo galinacea. - L.\cel{1}\sou{lira}. - DEFIRS.}
\l{}
\l{\sou{lirika.}\cel{2}I, Qua kontenas poemi drama, akompanata da kanti}
\l{\cel{3}e muziko. - II.Qua konservis (quankam ne destinita a}
\l{\cel{3}kanteso) la formo di la poezio lirika di la antiqui.}
\l{\cel{3}- DEFIRS.}
}
